\section{Problem Setup}
\label{sec:setup}

We study footstep planning for a Unitree Go1 quadruped walking forward
over ground with voids. At a fixed rate, a planner decides whether to
start a step and which one, and a low-level controller executes it.

At each tick, the planner observes the
robot's state, a set
$\mathcal C$ of candidate footholds, and a short history of the
robot's motion. It returns either the no-step option or a step
$(\ell, \mathbf u^\star, T)$: a leg $\ell$, a landing point
$\mathbf u^\star \in \mathcal C$, and a swing time $T$. The landing
point is an offset from the leg's hip in the yaw frame, the
gravity-aligned frame that turns with the robot's heading, taken at
the moment of the decision.

To execute the step, a controller $c$ receives a command $\mathbf u$,
also an offset from the hip in this yaw frame, and a start time. Sent
$\mathbf u$ for leg $\ell$ and swing time $T$ in robot state
$\mathbf s$, it lands the foot at
\begin{equation}
  \mathbf o = F_c(\mathbf s, \ell, T, \mathbf u)
    = \mathbf u + \boldsymbol\delta_c(\mathbf s, \ell, T, \mathbf u),
  \label{eq:forward}
\end{equation}
where $\boldsymbol\delta_c$ is the controller's landing error. Both
$\boldsymbol\delta_c$ and the controller's step timing differ between
controllers, and in general $\boldsymbol\delta_c \neq \mathbf 0$
(Sec.~\ref{subsec:method-idm}). An execution layer therefore sits
between the planner and the controller and turns each step into a
command and a start time for controller $c$.

We train the planner in
Isaac Lab with the \emph{training controller}, the convex MPC of Di
Carlo et al.\ \cite{dicarlo_dynamic_2018}, to follow a
forward velocity command over terrain with voids without falling
(Sec.~\ref{subsec:method-training}). We then freeze the planner's
weights and run it with OCS2
(Sec.~\ref{subsec:method-deployment}). Only the execution layer is
refit for the new controller. The goal is success on held-out
worlds comparable to that of OCS2's own gait.
